% !TeX root = ../main.tex
% Eigene Hilfsbefehle und mathematische Kurzschreibweisen.
% Diese Datei wird bereits in main.tex eingebunden.

\newcommand{\op}[1]{\operatorname{#1}}
\newcommand{\arr}[2]{\begin{array}{#1}#2\end{array}}
\renewcommand{\d}{\ensuremath{\text{d}}}

\newcommand{\NN}{\mathbb{N}}
\newcommand{\RR}{\mathbb{R}}
\newcommand{\CC}{\mathbb{C}}
\newcommand{\pdb}[2]{\frac{\partial #1}{\partial #2}}

% Markierung offener Stellen im Protokoll: \todo{Messwert prüfen}
\newcommand{\margtodo}{\marginpar{\textbf{\textcolor{expcolor}{ToDo}}}}
\newcommand{\todo}[1]{{\textbf{\textcolor{expcolor}{[\margtodo{}#1]}}}}
\newenvironment{deprecated}{\begin{color}{gray}}{\end{color}}

% Abbildungs-/Tabellenhilfen aus der ursprünglichen Vorlage
\newcommand{\xcaption}[2]{\caption[#1]{\textbf{#1} #2}}
\newcommand{\xfigure}[5]{%
    \begin{figure}[#1]
        \centering
        \includegraphics[scale=#2]{./fig/#3}
        \xcaption{#4}{#5}
        \label{fig:#3}
    \end{figure}}
\newcommand{\xfigurerot}[5]{%
    \begin{figure}[#1]
        \centering
        \includegraphics[angle=270,scale=#2]{./fig/#3}
        \xcaption{#4}{#5}
        \label{fig:#3}
    \end{figure}}
\newcommand{\xtable}[4]{%
    \begin{table}[#1]
        \centering
        \xcaption{#3}{#4}
        \rowcolors{3}{gray!10}{white}
        \include{./tab/#2}
        \label{tab:#2}
    \end{table}}

% Quantenmechanik und Mathematik
\newcommand{\ket}[1]{\left|#1\right\rangle}
\newcommand{\bra}[1]{\left\langle#1\right|}
\newcommand{\braket}[2]{\left\langle#1 \middle| #2\right\rangle}
\newcommand{\bratenket}[3]{\left\langle#1 \middle|\middle| #2 \middle|\middle| #3\right\rangle}
\newcommand{\anglemean}[1]{\left\langle #1 \right\rangle}
\newcommand{\norm}[1]{\left\lVert #1 \right\rVert}

\newcommand{\updownarrows}{\text{\rotatebox[origin=c]{90}{$\rightleftarrows$}}}
\newcommand{\downuparrows}{\text{\rotatebox[origin=c]{270}{$\rightleftarrows$}}}
\newcommand{\neswarrows}{\text{\rotatebox[origin=c]{45}{$\rightleftarrows$}}}
\newcommand{\swnearrows}{\text{\rotatebox[origin=c]{225}{$\rightleftarrows$}}}

\newcommand{\cre}{c^\dagger}
\newcommand{\anh}{c^{\vphantom{\dagger}}}
\newcommand{\numb}{n^{\vphantom{\dagger}}}
\newcommand{\fullstop}{\text{\,.}}
\newcommand{\comma}{\text{\,,}}
